\section{ОРГАНІЗАЦІЯ ТА МЕТОДИКА ПРОВЕДЕННЯ ДОСЛІДНО-ЕКСПЕРИМЕНТАЛЬНОЇ РОБОТИ}\label{chap3}

У попередніх розділах дисертації було здійснено теоретичний аналіз проблеми навчання розробки імерсивних освітніх ресурсів майбутніх учителів інформатики (розділ~1), визначено дидактичні підходи та принципи професійної підготовки, побудовано структурно-функціональну модель та обґрунтовано педагогічні умови навчання (розділ~2). Третій розділ присвячено організації та проведенню дослідно-експериментальної роботи з перевірки ефективності розробленої методики навчання розробки імерсивних освітніх ресурсів майбутніх учителів інформатики. Розроблено критеріальний апарат для оцінювання сформованості відповідних умінь, описано організацію та методику проведення педагогічного експерименту на базі Криворізького державного педагогічного університету (КДПУ), а також подано кількісний і якісний аналіз отриманих результатів із залученням методів математичної статистики.

Дослідно-експериментальна робота була спрямована на перевірку гіпотези дослідження, сформульованої у вступі: рівень сформованості вмінь майбутніх учителів інформатики у сфері розробки імерсивних освітніх ресурсів суттєво підвищиться за умов впровадження розробленої методики навчання, що спирається на структурно-функціональну модель підготовки та реалізацію визначених педагогічних умов.

\subsection{Критерії, показники та рівні сформованості вмінь майбутніх учителів інформатики у сфері розробки імерсивних освітніх ресурсів}\label{sec:3.1}

Оцінювання ефективності будь-якої педагогічної методики потребує розробки адекватного критеріального апарату, що відповідає меті та завданням дослідження. Для об'єктивного оцінювання ефективності запропонованої методики навчання розробки імерсивних освітніх ресурсів (ІОР) необхідно визначити критерії, показники та рівні сформованості відповідних умінь, а також розробити діагностичний інструментарій.

Спираючись на теоретичні положення, визначені у розділах~1--2, зокрема на структурно-функціональну модель підготовки (підрозділ~\ref{subsec:2.2}) та педагогічні умови навчання (підрозділ~\ref{subsec:2.3}), а також ураховуючи специфіку діяльності з проєктування та створення імерсивних освітніх ресурсів на основі технологій WebAR/VR/XR, було сформовано критеріальний апарат, що складається з чотирьох критеріїв, відповідних показників та чотирьох рівнів сформованості.

\subsubsection{Обґрунтування вибору критеріїв оцінювання}

Критерій (від грец. $\varkappa\rho\iota\tau\acute{\eta}\rho\iota o\nu$~-- мірило, засіб судження)~-- ознака, на підставі якої здійснюється оцінювання, визначення або класифікація будь-чого~\cite{Goncharenko1997Dictionary}. У педагогічних дослідженнях критерій виступає як об'єктивна мірка, що дає змогу виявити рівень досягнення поставлених навчальних цілей та ступінь сформованості компетентностей~\cite{Goncharenko2011Pedagogy}.

Вибір критеріїв оцінювання здійснювався на основі таких методологічних засад:
\begin{itemize}
\item відповідність компонентній структурі готовності, визначеній у структурно-функціональній моделі підготовки (підрозділ~\ref{subsec:2.2});
\item узгодженість із педагогічними умовами навчання (підрозділ~\ref{subsec:2.3});
\item відображення специфіки діяльності з розробки ІОР як технічно складної, творчої та педагогічно зумовленої;
\item забезпечення повноти та всебічності оцінювання;
\item вимірюваність та діагностичність кожного критерію.
\end{itemize}

На основі зазначених засад виділено чотири критерії: мотиваційний, когнітивний, операційно-діяльнісний та рефлексивно-оцінювальний. Ці критерії утворюють цілісну систему, що відповідає основним блокам структурно-функціональної моделі підготовки: мотиваційно-цільовому, теоретико-методологічному та змістовому, організаційно-технологічному та рефлексивно-оцінювальному.

\textbf{Мотиваційний критерій} характеризує ставлення майбутнього вчителя інформатики до процесу вивчення та розробки ІОР, його внутрішню мотивацію до опанування технологій WebAR/VR/XR, усвідомлення значущості імерсивних технологій у сучасному освітньому процесі. Обґрунтування включення мотиваційного критерію спирається на положення про визначальну роль мотивації у процесі навчання~\cite{Ziaziun2000Philosophy}. Розробка імерсивних освітніх ресурсів потребує значних інтелектуальних зусиль у засвоєнні нових технологій, зокрема бібліотек MindAR, Three.js, TensorFlow.js, A-Frame, стандарту WebXR Device API, що вимагає стійкої позитивної мотивації. Без достатнього рівня мотивації студенти не зможуть подолати труднощі, пов'язані з опануванням складного технологічного стеку, який включає програмування (JavaScript, WebGL), роботу з тривимірною графікою, комп'ютерним зором та машинним навчанням. Мотиваційний критерій відповідає мотиваційно-цільовому блоку структурно-функціональної моделі та забезпечується реалізацією психолого-педагогічних умов навчання.

\textbf{Когнітивний критерій} відображає рівень теоретичних знань майбутнього вчителя інформатики, необхідних для усвідомленої розробки ІОР. Розуміння принципів функціонування технологій WebAR~-- від низькорівневого WebGL до високорівневих бібліотек-обгорток (MindAR, A-Frame)~-- є необхідною умовою для ефективної практичної діяльності. Когнітивний критерій охоплює: знання сутності, видів та особливостей ІОР; розуміння архітектури сучасних бібліотек та фреймворків для розробки WebAR додатків; знання принципів тривимірного рендерингу у веб-браузері; розуміння алгоритмів комп'ютерного зору, що лежать в основі відстеження об'єктів; знання основ інтеграції моделей машинного навчання; а також знання методики використання ІОР у навчальному процесі. Цей критерій відповідає теоретико-методологічному та змістовому блокам моделі і забезпечується реалізацією методичних умов навчання.

\textbf{Операційно-діяльнісний критерій} характеризує рівень практичних умінь та навичок щодо безпосередньої розробки ІОР~-- від налаштування середовища розробки до створення повнофункціональних WebAR додатків із різними видами відстеження (зображень, облич, довкілля), інтеграцією моделей машинного навчання, застосуванням тривимірного моделювання та анімації. Цей критерій є центральним у системі оцінювання, оскільки саме практичні вміння розробки ІОР є кінцевою метою навчання, яке має яскраво виражений діяльнісний характер. Операційно-діяльнісний критерій відповідає організаційно-технологічному блоку моделі та забезпечується реалізацією технологічних умов навчання.

\textbf{Рефлексивно-оцінювальний критерій} відображає здатність майбутнього вчителя інформатики до самоаналізу та самооцінки власної діяльності з розробки ІОР, уміння аналізувати якість створених освітніх ресурсів з педагогічної та технічної точок зору, здатність до критичного оцінювання доцільності застосування імерсивних технологій у конкретних навчальних ситуаціях, а також готовність до подальшого самовдосконалення у цій сфері. Рефлексія є невід'ємним компонентом професійного розвитку вчителя~\cite{ReflexiveApproach2018}, адже здатність аналізувати власний досвід та коригувати діяльність на основі цього аналізу є запорукою безперервного професійного зростання. Цей критерій відповідає рефлексивно-оцінювальному блоку моделі та забезпечується реалізацією організаційних умов навчання.

\subsubsection{Показники за кожним критерієм}

Для кожного з визначених критеріїв було виділено конкретні показники, що дають змогу оцінити рівень сформованості відповідних компетентностей і є основою для розробки діагностичного інструментарію.

\textit{Показники мотиваційного критерію:}
\begin{enumerate}
\item стійкість інтересу до технологій доповненої та віртуальної реальності та їх освітніх застосувань~-- визначається тривалістю та систематичністю залученості студента до вивчення WebAR/VR/XR технологій;
\item усвідомлення значущості ІОР для підвищення якості навчального процесу~-- виражається у здатності аргументувати переваги використання імерсивних ресурсів у шкільному курсі інформатики та інших дисциплінах;
\item прагнення до самостійного опанування нових засобів розробки WebAR/VR/XR~-- проявляється у ініціативному вивченні нових бібліотек та технологій, що виходять за межі обов'язкової програми;
\item орієнтація на використання імерсивних технологій у майбутній професійній діяльності~-- виражається у планах щодо впровадження ІОР у власну педагогічну практику;
\item готовність до подолання труднощів при освоєнні складних технологій~-- визначається наполегливістю студента у вирішенні технічно складних завдань, що потребують знання JavaScript, Three.js, WebGL, TensorFlow.js;
\item мотивація до створення власних імерсивних освітніх ресурсів~-- проявляється у розробці ІОР з власної ініціативи для конкретних предметних галузей.
\end{enumerate}

\textit{Показники когнітивного критерію:}
\begin{enumerate}
\item знання сутності, видів та класифікації ІОР~-- включає розуміння відмінностей між доповненою реальністю (AR), віртуальною реальністю (VR), змішаною реальністю (MR) та розширеною реальністю (XR), а також специфіки веб-орієнтованих рішень;
\item розуміння принципів функціонування WebAR~-- охоплює знання процесу отримання відеопотоку з камери пристрою, тривимірного рендерингу засобами WebGL та Three.js, алгоритмів відстеження об'єктів (маркерне, безмаркерне, відстеження природних зображень~-- NFT);
\item знання архітектури та можливостей основних бібліотек~-- включає розуміння структури MindAR (модулі відстеження зображень та облич), Three.js (сцена, камера, рендерер, об'єкти, матеріали, освітлення, анімація), A-Frame (декларативний підхід до побудови WebAR сцен);
\item розуміння стандарту WebXR Device API~-- охоплює знання принципів відстеження довкілля (world tracking), механізму hit testing для розміщення об'єктів на реальних поверхнях, управління контролерами XR-пристроїв;
\item знання основ інтеграції моделей машинного навчання у WebAR~-- включає розуміння архітектури TensorFlow.js, застосування попередньо навчених моделей (handpose для відстеження рухів рук, face-api.js для розпізнавання емоцій), принципів навчання власних моделей за допомогою Google Teachable Machine;
\item розуміння форматів тривимірних моделей та принципів їх використання~-- включає знання форматів GLTF та GLB, принципів оптимізації 3D-моделей для веб-застосунків, роботи з текстурами та матеріалами;
\item знання методики проєктування та використання ІОР у навчальному процесі~-- включає розуміння принципів педагогічного дизайну імерсивних ресурсів, критеріїв доцільності їх застосування;
\item обізнаність щодо сучасних комерційних та некомерційних WebAR платформ~-- включає знання можливостей та обмежень 8th Wall, Zappar, model-viewer та їх порівняння з MindAR.
\end{enumerate}

\textit{Показники операційно-діяльнісного критерію:}
\begin{enumerate}
\item уміння налаштовувати середовище розробки WebAR~-- включає встановлення та налаштування локального сервера (live-server, http-server), використання ngrok для тунелювання HTTPS-з'єднань, налаштування віддаленого налагодження на мобільних пристроях через Chrome DevTools;
\item здатність створювати тривимірні сцени засобами Three.js~-- охоплює створення об'єктів з різною геометрією (BoxGeometry, SphereGeometry, ConeGeometry, PlaneGeometry, TorusGeometry), застосування матеріалів (MeshBasicMaterial, MeshStandardMaterial, MeshPhongMaterial) та текстур, налаштування освітлення (AmbientLight, DirectionalLight, PointLight, SpotLight), управління камерами;
\item уміння розробляти WebAR додатки з відстеженням зображень~-- включає підготовку зображень для відстеження, компіляцію NFT маркерів за допомогою MindAR Image Compiler, створення прив'язок (anchors) тривимірних об'єктів до відстежуваних зображень, завантаження відео, аудіо та 3D-моделей як AR-контенту;
\item уміння розробляти WebAR додатки з відстеженням обличчя~-- охоплює використання модуля MindAR Face Tracking, роботу з опорними точками обличчя (face landmarks), створення окклюдерів для реалістичного накладання віртуальних об'єктів, застосування масок та фільтрів, реалізацію перемикання між передньою та задньою камерами пристрою;
\item уміння розробляти WebAR додатки з відстеженням довкілля~-- включає використання WebXR Device API, реалізацію механізму hit testing для визначення поверхонь реального світу, розміщення віртуальних об'єктів на цих поверхнях, управління контролерами для взаємодії з AR-вмістом;
\item здатність інтегрувати моделі машинного навчання у WebAR~-- охоплює застосування трикрокової методики: інтеграцію стандартних моделей TensorFlow.js (handpose для розпізнавання жестів рук), навчання власних моделей класифікації зображень за допомогою Teachable Machine та їх експорт до TensorFlow.js, інтеграцію навчених моделей у повнофункціональні WebAR додатки;
\item уміння завантажувати та використовувати зовнішні тривимірні моделі у форматах GLTF/GLB за допомогою GLTFLoader у Three.js;
\item здатність розробляти навчальні AR-додатки для конкретних предметних галузей, зокрема на прикладі віртуальної хімічної лабораторії у доповненій реальності;
\item уміння використовувати A-Frame та model-viewer для швидкого прототипування WebAR сцен;
\item навички роботи з мультимедійним контентом (відео, аудіо) у WebAR додатках.
\end{enumerate}

\textit{Показники рефлексивно-оцінювального критерію:}
\begin{enumerate}
\item здатність до самоаналізу процесу та результатів власної діяльності з розробки ІОР~-- виражається в умінні виявити власні сильні та слабкі сторони у процесі розробки;
\item уміння критично оцінювати якість створених імерсивних ресурсів~-- включає оцінку за технічними (продуктивність, стабільність, крос-платформність), юзабіліті (зручність використання, інтуїтивність інтерфейсу) та педагогічними (відповідність навчальним цілям, ефективність засвоєння матеріалу) параметрами;
\item здатність аналізувати доцільність обрання конкретних технологій WebAR~-- виражається в умінні обґрунтувати вибір між MindAR, A-Frame, WebXR, 8th Wall, Zappar для розв'язання конкретних освітніх завдань;
\item уміння формулювати рекомендації щодо вдосконалення розроблених ІОР~-- включає здатність визначити напрями покращення та запропонувати конкретні технічні й методичні рішення;
\item готовність до взаємооцінювання результатів роботи колег~-- виражається в умінні надавати конструктивний зворотний зв'язок, формулювати об'єктивні зауваження та пропозиції;
\item здатність планувати подальший саморозвиток у сфері імерсивних технологій~-- включає визначення власних освітніх потреб та планування шляхів їх задоволення;
\item уміння оцінювати педагогічну ефективність застосування ІОР у навчальному процесі~-- включає аналіз впливу використання імерсивних ресурсів на мотивацію та результати навчання учнів.
\end{enumerate}

\subsubsection{Характеристика рівнів сформованості вмінь}

На основі визначених критеріїв та показників було виділено чотири рівні сформованості вмінь майбутніх учителів інформатики у сфері розробки ІОР: низький, середній, достатній та високий. Такий розподіл є загальноприйнятим у педагогічних дослідженнях і відповідає чотирирівневій системі оцінювання~\cite{Fitsula2006Pedagogics}.

\textit{Низький рівень} характеризується нестійкою або відсутньою мотивацією до вивчення технологій доповненої та віртуальної реальності; фрагментарними знаннями про сутність ІОР та принципи роботи WebAR; здатністю виконувати лише найпростіші дії з розробки WebAR додатків за покроковою інструкцією без розуміння їх суті; відсутністю здатності до самостійного аналізу та оцінювання результатів власної діяльності. Студент із низьким рівнем здатний відтворити базові приклади з навчального посібника (наприклад, створити просту сцену Three.js з одним об'єктом або прив'язати зображення до маркера MindAR), але не може самостійно модифікувати або розширювати їх для вирішення нових завдань.

\textit{Середній рівень} характеризується ситуативною мотивацією, що залежить від зовнішніх факторів; базовими знаннями основних понять та принципів роботи WebAR технологій, але з прогалинами у розумінні окремих компонентів; здатністю розробляти прості WebAR додатки з одним типом відстеження за зразком з незначними модифікаціями; здатністю до часткового самоаналізу результатів роботи за заданими критеріями. Студент із середнім рівнем здатний створити AR-візитівку або простий додаток з накладанням текстури та 3D-об'єкта на маркер, але стикається зі значними труднощами при вирішенні нестандартних завдань.

\textit{Достатній рівень} характеризується стійкою внутрішньою мотивацією та усвідомленням значущості імерсивних технологій в освіті; системними знаннями основних технологій розробки WebAR (MindAR, Three.js, WebXR) та розумінням їх взаємозв'язків; здатністю самостійно розробляти WebAR додатки різних типів (з відстеженням зображень, облич, довкілля) та інтегрувати базові моделі машинного навчання; здатністю до системного аналізу та оцінювання якості створених ІОР за різними параметрами. Студент із достатнім рівнем здатний розробити повноцінний навчальний WebAR додаток, але може потребувати консультації при вирішенні технічно складних завдань.

\textit{Високий рівень} характеризується стійкою внутрішньою мотивацією та прагненням до інноваційної діяльності у сфері ІОР; глибокими та системними знаннями всього спектру технологій WebAR/VR/XR, включаючи комерційні SDK; здатністю самостійно проєктувати та реалізовувати складні імерсивні освітні ресурси з використанням кількох типів відстеження, інтеграцією власних моделей машинного навчання, складною тривимірною графікою та анімацією; здатністю до творчого підходу в оцінюванні, прогнозування результатів та формулювання обґрунтованих рекомендацій. Студент із високим рівнем здатний самостійно проєктувати та реалізовувати такі складні ІОР, як віртуальна хімічна лабораторія у доповненій реальності, інтегрувати розпізнавання жестів рук для керування AR-контентом, навчити власну модель класифікації зображень та вбудувати її у WebAR додаток.

У табл.~\ref{tab:criteria_levels_ier} подано узагальнену характеристику рівнів сформованості вмінь за кожним із визначених критеріїв.

\begin{table}[!h]
\centering
\caption{Характеристика рівнів сформованості вмінь за критеріями}
\label{tab:criteria_levels_ier}
\small
\begin{tabularx}{\linewidth}{|p{2cm}|p{3.2cm}|p{3cm}|p{3.4cm}|X|}
\hline
\hfil \textbf{Критерій} & \hfil \textbf{Низький} & \hfil \textbf{Середній} & \hfil \textbf{Достатній} & \hfil \textbf{Високий} \\
\hline
Мотива\-ційний &
Нестійка мотивація; відсутність інтересу до імерсивних технологій; пасивне ставлення до навчальної діяльності &
Ситуативна мотивація; інтерес залежить від зовнішніх факторів; епізодична активність &
Стійка внутрішня мотивація; усвідомлення значущості ІОР в освіті; систематична активність у навчанні &
Стійка внутрішня мотивація; прагнення до інноваційної діяльності; самостійне вивчення нових технологій WebAR/VR/XR \\
\hline
Когні\-тивний &
Фрагментарні знання; нерозуміння базових принципів WebAR; плутання понять AR, VR, MR, XR &
Базові знання з прогалинами; розуміння окремих компонентів без цілісної картини; знання основ JavaScript та HTML &
Системні знання; розуміння архітектури MindAR, Three.js, WebXR; знання принципів інтеграції TensorFlow.js &
Глибокі системні знання; розуміння оптимізації та алгоритмів комп'ютерного зору; знання комерційних SDK \\
\hline
Опера\-ційно-діяль\-нісний &
Виконання дій за покроковою інструкцією; відтворення базових прикладів без модифікацій &
Розробка простих додатків з одним типом відстеження; незначні модифікації зразків &
Самостійна розробка WebAR додатків різних типів; інтеграція базових ML моделей; розробка ІОР з педагогічним сценарієм &
Проєктування складних ІОР з кількома типами відстеження; інтеграція власних ML моделей; створення навчальних AR-лабораторій \\
\hline
Рефлек\-сивно-оціню\-вальний &
Відсутність здатності до самоаналізу; неспроможність оцінити якість ІОР &
Частковий самоаналіз за заданими критеріями; оцінювання з допомогою викладача &
Системний самоаналіз; критичне оцінювання якості ІОР за технічними та педагогічними параметрами &
Творчий підхід до оцінювання; прогнозування результатів впровадження ІОР; організація взаємо\-оцінювання \\
\hline
\end{tabularx}
\end{table}

\subsubsection{Оцінювальна шкала та рубрика}

Для кількісного оцінювання рівня сформованості вмінь було розроблено оцінювальну шкалу, що передбачає комплексне оцінювання за всіма чотирма критеріями. Кожен критерій оцінювався за 25-бальною шкалою, що в сумі давало максимум 100~балів. Такий підхід забезпечує рівномірну вагу кожного критерію у загальній оцінці та відповідає 100-бальній системі оцінювання, прийнятій у КДПУ.

Оцінка за кожним критерієм визначалася на основі комплексу діагностичних процедур, що включали тестування, оцінювання практичних завдань, експертне оцінювання та самооцінювання. Остаточний бал за критерієм розраховувався як зважена сума оцінок за окремими показниками.

У табл.~\ref{tab:assessment_rubric} подано рубрику оцінювання, що визначає відповідність балів рівням сформованості.

\begin{table}[!h]
\centering
\caption{Рубрика оцінювання рівнів сформованості вмінь розробки ІОР}
\label{tab:assessment_rubric}
\small
\begin{tabular}{|l|c|c|c|c|}
\hline
\hfill \textbf{Критерій} & \textbf{Низький} & \textbf{Середній} & \textbf{Достатній} & \textbf{Високий} \\
 & (0--6 б.) & (7--13 б.) & (14--19 б.) & (20--25 б.) \\
\hline
Мотиваційний (до 25~б.) & 0--6 & 7--13 & 14--19 & 20--25 \\
\hline
Когнітивний (до 25~б.) & 0--6 & 7--13 & 14--19 & 20--25 \\
\hline
Операційно-діяльнісний (до 25~б.) & 0--6 & 7--13 & 14--19 & 20--25 \\
\hline
Рефлексивно-оцінювальний (до 25~б.) & 0--6 & 7--13 & 14--19 & 20--25 \\
\hline
\textbf{Загальний бал (до 100~б.)} & \textbf{0--24} & \textbf{25--52} & \textbf{53--76} & \textbf{77--100} \\
\hline
\end{tabular}
\end{table}

\subsubsection{Інструменти діагностики}

Для діагностики рівня сформованості вмінь за кожним критерієм було розроблено та апробовано комплекс діагностичних інструментів.

\textit{Діагностика мотиваційного критерію} здійснювалася за допомогою:
\begin{itemize}
\item авторської анкети <<Мотивація до вивчення імерсивних технологій>> (25 питань за шкалою Лікерта), що включала блоки: внутрішня мотивація, зовнішня мотивація, ціннісне ставлення до ІОР, професійна спрямованість;
\item систематичного спостереження за активністю студентів під час навчальних занять та лабораторних робіт (участь у дискусіях, ініціативність, завершеність завдань);
\item аналізу самостійної роботи студентів у системі Moodle курсу <<Синтетична реальність в освіті>>~-- регулярність виконання завдань (18 тижневих модулів), своєчасність здачі робіт, участь у форумах;
\item бесід зі студентами щодо їхніх планів використання імерсивних технологій у педагогічній практиці.
\end{itemize}

\textit{Діагностика когнітивного критерію} здійснювалася за допомогою:
\begin{itemize}
\item комплексу тестових завдань (40 питань) з різних розділів навчального посібника, що охоплюють: основи JavaScript та HTML для WebAR, принципи 3D-рендерингу та Three.js, відстеження зображень за допомогою MindAR, відстеження облич, WebXR Device API та відстеження довкілля, інтеграцію TensorFlow.js, A-Frame та інші технології WebAR, методику використання ІОР в освіті;
\item контрольних запитань до кожного з 10 розділів навчального посібника, що потребували розгорнутих відповідей з поясненням принципів роботи технологій;
\item експертного оцінювання відповідей на проблемні запитання, що потребують аналізу, порівняння та узагальнення.
\end{itemize}

\textit{Діагностика операційно-діяльнісного критерію} здійснювалася за допомогою:
\begin{itemize}
\item оцінювання результатів виконання лабораторних робіт: лабораторна робота~1 <<Візитівка у доповненій реальності>> (створення WebAR-візитівки з відстеженням зображень засобами MindAR), лабораторна робота~2 <<Примірка віртуальних аксесуарів>> (розробка додатку з відстеженням обличчя засобами MindAR Face Tracking з використанням окклюдерів);
\item оцінювання комплексу із 10 варіантів творчих завдань з 3D-моделювання та AR: побудова масштабованих моделей Землі та Місяця, пірамід Гізи, Ханойських веж, Сатурна з кільцями, планетарної моделі атома літію, будівлі Криворізької міської ради, фізичних симуляцій тощо з прив'язкою моделей до маркерів доповненої реальності;
\item оцінювання індивідуальних підсумкових проєктів~-- розробка повноцінного навчального WebAR ресурсу для конкретної предметної галузі з педагогічним обґрунтуванням;
\item аналізу портфоліо студентських робіт;
\item оцінювання якості реалізації віртуальної хімічної лабораторії в доповненій реальності (anchem-ar)~-- інтеграційного завдання, що потребує одночасного використання кількох маркерів, обробки їх взаємного розташування та відтворення результатів хімічних реакцій.
\end{itemize}

\textit{Діагностика рефлексивно-оцінювального критерію} здійснювалася за допомогою:
\begin{itemize}
\item есе-рефлексій студентів щодо досвіду розробки ІОР, в яких вони аналізували труднощі, здобутки та перспективи використання набутих умінь;
\item захисту індивідуальних підсумкових проєктів з обґрунтуванням обраних технічних рішень, оцінкою їх ефективності та визначенням напрямів удосконалення;
\item процедури взаємооцінювання студентських проєктів, під час якої кожен студент оцінював роботи щонайменше двох колег за визначеними критеріями (технічна якість, педагогічна доцільність, оригінальність);
\item підсумкового самооцінювання рівня готовності до розробки ІОР за допомогою спеціально розробленого опитувальника (20 питань за шкалою Лікерта).
\end{itemize}

\subsubsection{Зв'язок критеріїв з компонентами структурно-функціональної моделі та педагогічними умовами}

Визначені критерії оцінювання утворюють цілісну систему, безпосередньо пов'язану з блоками структурно-функціональної моделі підготовки (підрозділ~\ref{subsec:2.2}) та педагогічними умовами навчання (підрозділ~\ref{subsec:2.3}). Цей зв'язок забезпечує внутрішню узгодженість дослідження: критерії оцінюють саме ті компетентності, формування яких передбачено моделлю, а педагогічні умови забезпечують середовище, в якому це формування відбувається.

У табл.~\ref{tab:criteria_model_link} подано відповідність між критеріями оцінювання, блоками структурно-функціональної моделі, педагогічними умовами та основними інструментами діагностики.

\begin{table}[!h]
\centering
\caption{Відповідність критеріїв оцінювання блокам моделі та педагогічним умовам}
\label{tab:criteria_model_link}
\small
\begin{tabularx}{\linewidth}{|p{2.5cm}|p{3cm}|p{2.7cm}|X|c|}
\hline
\hfil \textbf{Критерій} & \hfil \textbf{Блок моделі} & \textbf{\makecell{Педагогічна \\умова}} & \textbf{\makecell{Основні інстру\-менти \\діаг\-ностики}} & \textbf{\makecell{Максимальний \\бал}} \\
\hline
Мотива\-ційний & Мотиваційно-цільовий & Психолого-педагогічна & Анкетування, спостереження, аналіз активності у Moodle & 25 \\
\hline
Когнітив\-ний & Теоретико-методо\-логічний, змістовий & Методична & Тестування (40 питань), контрольні запитання, експертна оцінка & 25 \\
\hline
Операційно-діяль\-нісний & Організаційно-техно\-логічний & Технологічна & Лабораторні роботи, творчі завдання, проєкти, портфоліо & 25 \\
\hline
Рефлексивно-оціню\-вальний & Рефлексивно-оціню\-вальний & Організаційна & Есе-рефлексії, захист проєктів, взаємо\-оцінювання & 25 \\
\hline
\end{tabularx}
\end{table}

Таким чином, розроблений критеріальний апарат для оцінювання рівня сформованості вмінь майбутніх учителів інформатики у сфері розробки ІОР включає: чотири критерії (мотиваційний, когнітивний, операційно-діяльнісний та рефлексивно-оцінювальний), що відповідають блокам структурно-функціональної моделі підготовки; конкретні показники для кожного критерію; чотири рівні сформованості (низький, середній, достатній, високий) із детальною характеристикою кожного рівня за кожним критерієм; 100-бальну оцінювальну шкалу з рівномірним розподілом між критеріями; а також комплекс діагностичних інструментів. Розроблений апарат було використано для оцінювання результатів педагогічного експерименту, організацію та методику проведення якого описано у наступному підрозділі.
